\honest{Hindsight Devalues Science}{Eliezer Yudkowsky}{2007}

This essay is closely based on an excerpt from a social psychology textbook.

An editor of \textsc{The Atlantic} said that social scientists discover every day that ``people's behavior is pretty much what you'd expect.'' That expectation ``is all hindsight.'' \nb{The evidence the post gives shows that the feeling of having expected a finding is unreliable. It does not show that no finding could have been expected.}

The historian Arthur Schlesinger, Jr. called studies of soldiers in the Second World War ``ponderous demonstrations'' of common sense. I give five of their findings, each with a plausible explanation: better-educated soldiers adjusted worse, Southern soldiers coped better with the tropical heat, and so on. How many could you have predicted?

All five are the opposite of what was found. \nb{The list comes from Paul Lazarsfeld's 1949 review, where it had six items. For the climate item, his own summary of the true finding is that Southerners ``showed no greater ability than Northerners'': no difference, not a reversal.} ``That's how good your model really was. The measure of your strength as a rationalist is your ability to be more confused by fiction than by reality.'' \nb{This is the thesis of ``Your Strength as a Rationalist,'' posted six days earlier, and it has the problem noted there. Lazarsfeld drew a different lesson from his list: ``Since every kind of human reaction is conceivable,'' only data can tell which occurs. Without the data, a reader has no ground to be more confused by one list than the other.}

Unless I reversed the results again. Does your thinking now, when you really do not know, feel different from rationalizing a known answer?

Daphna Baratz gave students pairs of findings, one true and one its opposite, and they rated each as what they ``would have predicted.'' This is why people think they do not need science.

Hindsight makes us undervalue how surprising findings are, ``especially the discoveries we understand,'' and if you know neurology or physics you ``probably'' underestimate the surprise there too. I give no evidence for this. \nb{Studies of experts and hindsight bias disagree: one meta-analysis found experts less prone to it, a later and larger one found no relation.}

My remedy: ``We need to make a conscious effort to be shocked enough.'' \nb{Research on hindsight bias suggests that knowing about it and intending to avoid it do not remove it. What reduces it is explaining how the outcome that did not happen could have happened. The post's question about rationalizing either side comes close to that; its last line does not.}
